\chapter{Method}

This chapter describes the methods used to optimize the Majorana modes in the junction setup and to analyze their properties. We first discuss the parameter optimization process, including the model parameterization, loss function design based on Majorana criteria, and the optimization strategy. Next, we detail the post-optimization analysis of the setup, covering Hamiltonian reconstruction, parity classification, and various diagnostics. Finally, we describe the braiding protocol implemented in an effective Majorana model and how we validate the results using gate-validation checks.

\section{Numerical PMM Hamiltonian and Basis}
The numerical calculations are built from the spinless $n$-dot PMM Hamiltonian introduced in the theory chapter. This section specifies how that model is represented in the code: which Hamiltonian is diagonalized, which basis is used, how the fermionic operators are constructed, and how the model parameters are organized before optimization and braiding calculations.

\subsection{Single-Chain Hamiltonian}
For a single $n$-dot PMM chain, the Hamiltonian used in the optimization and diagnostic calculations is
\begin{equation}
H_{\mathrm{PMM}}
=
\sum_{i=1}^{n} \epsilon_i n_i
+ \sum_{i=1}^{n-1}
\left[
-t_i\left(c_i^\dagger c_{i+1}+c_{i+1}^\dagger c_i\right)
+\Delta_i\left(c_i c_{i+1}+c_{i+1}^\dagger c_i^\dagger\right)
+U_i n_i n_{i+1}
\right].
\label{eq:method_pmm_hamiltonian}
\end{equation}
Here $c_i^\dagger$ and $c_i$ create and annihilate a spinless fermion on dot $i$, and $n_i=c_i^\dagger c_i$ is the corresponding number operator. The onsite energies $\epsilon_i$ play the role of local chemical potentials, $t_i$ are elastic cotunneling amplitudes, $\Delta_i$ are crossed-Andreev pairing amplitudes, and $U_i$ are nearest-neighbor density-density interaction strengths.

\subsection{Many-Body Fock Basis}
All operators are represented in the full many-body Fock basis. For a chain with $n$ dots, a basis state is specified by the occupation pattern
\begin{equation}
\ket{n_1 n_2 \cdots n_n},
\qquad
n_i\in\{0,1\},
\end{equation}
so the Hilbert-space dimension of a single chain is $2^n$. Diagonalizing Eq.~\ref{eq:method_pmm_hamiltonian} in this basis gives the full many-body spectrum and eigenstates used to calculate the Majorana diagnostics.

For braiding calculations, several optimized chains are combined into a larger Hilbert space. If $D$ identical PMM subsystems are used, each containing $n$ dots, the total number of fermionic sites is $N_{\mathrm{tot}}=Dn$ and the full Hilbert-space dimension is $2^{N_{\mathrm{tot}}}$. In the three-arm braiding geometry used later, $D=3$, corresponding to subsystems $A$, $B$, and $C$.

\subsection{Jordan-Wigner Construction of Fermionic Operators}
\label{sec:method-jordan-wigner-construction}
Using the Jordan-Wigner construction introduced in Sec.~\ref{sec:theory-jordan-wigner}, the code builds fermionic creation and annihilation operators from local spin raising and lowering operators. For site $j$, the operators are represented schematically as
\begin{equation}
c_j^\dagger =
\left[\prod_{k<j}(-\sigma_k^z)\right]\sigma_j^+,
\qquad
c_j =
\left[\prod_{k<j}(-\sigma_k^z)\right]\sigma_j^-.
\label{eq:method_jw_operators}
\end{equation}
The same construction is used for single-chain calculations and for the enlarged three-subsystem Hilbert space.

Once the creation and annihilation operators have been constructed, the code precomputes the basic operator products needed repeatedly:
\begin{equation}
n_i=c_i^\dagger c_i,\qquad
h_{ij}=c_i^\dagger c_j+c_j^\dagger c_i,\qquad
p_{ij}=c_i^\dagger c_j^\dagger+c_j c_i,\qquad
d_{ij}=n_i n_j .
\end{equation}
These building blocks are then combined with the parameter values in Eq.~\ref{eq:method_pmm_hamiltonian} to assemble the Hamiltonian matrix. 
For the three-subsystem geometry, the Jordan-Wigner ordering is
$$
A_1, A_2, \ldots, A_n, B_1, B_2, \ldots, B_n, C_1, C_2, \ldots, C_n.
$$
Operators acting on subsystem B therefore carry a Jordan-Wigner string over all sites in subsystem A, while all operators acting on subsystem C carry a string over all sites in both A and B.

The parameter conventions used when assembling this Hamiltonian during optimization are described in Sec.~\ref{sec:method-optimization-vector}.

\section{Parameter Optimization and Loss Function}
The purpose of the optimization is to find parameter regimes where the finite PMM chain has the spectral and local signatures of Majorana-like states. The theory chapter defined the physical criteria; here they are turned into a numerical objective. For each trial parameter vector, the Hamiltonian in Eq.~\ref{eq:method_pmm_hamiltonian} is constructed, diagonalized, separated into parity sectors, and assigned a scalar loss. The optimizer then updates only the free parameters, while fixed parameters are reinserted when the physical Hamiltonian is rebuilt.

\subsection{Optimization Vector and Parameter Reconstruction}
\label{sec:method-optimization-vector}
The Hamiltonian in Eq.~\ref{eq:method_pmm_hamiltonian} is specified by the physical parameters $\{t_i,U_i,\epsilon_i,\Delta_i\}$. In the numerical implementation, each group of parameters can be fixed, homogeneous, or inhomogeneous. A fixed parameter is held at a prescribed value and is not included in the optimization vector. A homogeneous parameter has the same value on every site or bond, while an inhomogeneous parameter is allowed to vary independently from site to site or bond to bond. This makes it possible to compare restricted ansatz with more flexible parameter searches using the same Hamiltonian-building routine.

The optimizer therefore does not act directly on a full dictionary of physical parameters. Instead, it acts on a reduced real vector $\theta$ containing only the entries that are free in the chosen parameter configuration. For example, if $t_i$ is fixed and $U_i$ is fixed to a selected interaction strength, then those entries are excluded from $\theta$. If a parameter is homogeneous, only one number is included, or if it is inhomogeneous, one number is included for each site or bond.

The reduced vector is converted back into a parameter dictionary before the Hamiltonian is evaluated. Schematically,
\begin{equation}
\theta
\longrightarrow
\{t_i,U_i,\epsilon_i,\Delta_i\}
\longrightarrow
H_{\mathrm{PMM}}(\theta).
\end{equation}
During this conversion, fixed values are reinserted, homogeneous values are expanded across all sites or bonds when needed, and positive couplings such as $t_i$, $U_i$, and $\Delta_i$ are constrained to be non-negative. This separation lets the optimizer work with a compact vector while the Hamiltonian is still assembled from physical parameters.

In the optimization runs, the hopping amplitudes are fixed to $t_i=1$, while the interaction strength $U$ is fixed to the value being studied. The remaining parameters, the onsite energies $\epsilon_i$ and pairing amplitudes $\Delta_i$, are optimized. Fixing the hopping scale prevents the optimizer from lowering the loss through a trivial global rescaling of the kinetic energy, and instead forces the search to find nontrivial onsite and pairing profiles for each interaction strength. This also rules out the possibility of blowing up all other parameters in order to suppress the relative effect of interactions, which would be an unphysical way to achieve a low loss in the strongly interacting regime.

For each system size $n$ and interaction strength $U$, the lowest-loss physical parameter set is stored and later reused in the diagnostic and braiding calculations. The same parameter convention is also used when constructing the three-subsystem braiding geometry, as the optimized single-chain parameters define each subsystem, while additional inter-subsystem couplings are introduced only in the braiding-model sections below.

\subsection{Parity-Resolved Spectral Loss}
For each trial Hamiltonian, the full many-body spectrum is obtained by exact diagonalization. The eigenstates are then classified by the total fermion parity operator,
\begin{equation}
P_{\mathrm{tot}}=\prod_{i=1}^{n}(1-2n_i).
\end{equation}
This gives two ordered energy lists, one in the even sector and one in the odd sector,
\begin{equation}
\{E_k^{(e)}\},
\qquad
\{E_k^{(o)}\}.
\end{equation}
The loss function compares these parity partners level by level. For the parity-conserving Hamiltonians considered here, the even and odd sectors have equal dimension. The implementation nevertheless checks this explicitly, so that numerical or construction errors in the parity classification do not enter the loss evaluation.

The first loss contribution penalizes the energy splitting between even- and odd-parity partners,
\begin{equation}
\delta_k = \left|E_k^{(e)}-E_k^{(o)}\right|.
\end{equation}
As discussed in Sec.~\ref{sec:theory-parity-resolved-spectrum}, small even--odd splittings are one of the main spectral signatures of Majorana-like states. In the loss, these splittings are weighted more strongly at low energy. If $N_p$ parity pairs are used, the implementation defines a decreasing weight array
\begin{equation}
w_{\ell}
=
\left[
w_{\mathrm{max}}
-
\frac{\ell}{2N_p-2}
\left(
w_{\mathrm{max}}-w_{\mathrm{min}}
\right)
\right]^2,
\qquad
\ell=0,\ldots,2N_p-2,
\end{equation}
with $w_{\mathrm{max}}=3$ and $w_{\mathrm{min}}=0.1$. The degeneracy contribution is
\begin{equation}
\mathcal{L}_{\mathrm{deg}}
=
\sum_{k=0}^{N_p-1} w_k \left|E_k^{(e)}-E_k^{(o)}\right|.
\end{equation}
This term encourages the optimizer to produce an approximately parity-paired spectrum, with particular emphasis on the lowest-energy states.

The first $N_p$ entries of the weight array are used for the degeneracy term, while the remaining $N_p-1$ entries are used for the excitation-gap term. For adjacent parity pairs, the relevant gap is taken as
\begin{equation}
\Delta_k
=
\min\left(
E_{k+1}^{(e)}-E_k^{(e)},
E_{k+1}^{(o)}-E_k^{(o)}
\right).
\end{equation}
This $\Delta_k$ is the smaller spacing between adjacent levels in the two parity sectors. The corresponding weight is $w_{N_p+k}$.

Small gaps are penalized using a smooth threshold function,
\begin{equation}
\mathcal{L}_{\mathrm{gap}}
=
\sum_{k=0}^{N_p-2} w_{N_p+k}\,
\mathrm{softplus}
\left(
\Delta_{\mathrm{target}}-\Delta_k
\right),
\end{equation}
where $\Delta_{\mathrm{target}}$ is the desired minimum gap scale, and the softplus function is defined as $\mathrm{softplus}(x) = \log(1 + e^x)$\cite{NIPS2000_44968aec}.


\subsection{Charge and Majorana-Polarization Terms}
The local part of the loss turns the charge and Majorana-polarization diagnostics from Secs.~\ref{sec:theory-local-distinguishability} and \ref{sec:theory-majorana-polarization} into numerical penalties. The charge penalty measures how differently the even and odd partners occupy each site,
\begin{equation}
\mathcal{L}_{Q}
=
\sum_{k}\sum_{j}
\left|
\bra{e_k}n_j\ket{e_k}
-
\bra{o_k}n_j\ket{o_k}
\right|.
\end{equation}
This is the implementation form of the site-resolved charge-difference diagnostic.

The Majorana-polarization diagnostic itself was defined in Sec.~\ref{sec:theory-majorana-polarization}. In the optimizer, the corresponding penalty is implemented directly from the parity-changing matrix elements of the two local Majorana components $\gamma_{j,1}=c_j^\dagger+c_j$ and $\gamma_{j,2}=i(c_j^\dagger-c_j)$. For parity pair $k$ and site $j$, the code forms
\begin{equation}
\widetilde{M}_{k,j}
=
\left(
\bra{o_k}\gamma_{j,1}\ket{e_k}
\right)^2
+
\left(
\bra{o_k}\gamma_{j,2}\ket{e_k}
\right)^2 .
\end{equation}
The target profile is concentrated on the two ends of the chain with opposite sign,
\begin{equation}
\widetilde{M}^{\mathrm{tar}}_{k,1}=1,
\qquad
\widetilde{M}^{\mathrm{tar}}_{k,n}=-1,
\qquad
\widetilde{M}^{\mathrm{tar}}_{k,j}=0
\quad
(1<j<n).
\end{equation}
The sign of this target pattern is allowed to flip, since exchanging the labels of the two Majorana components changes the overall sign but not the physical localization pattern. The implemented Majorana-polarization penalty is therefore
\begin{equation}
\mathcal{L}_{\mathrm{MP}}
=
\min_{\sigma=\pm1}
\left|
\sum_{k,j}
\left(
\widetilde{M}_{k,j}
-
\sigma \widetilde{M}^{\mathrm{tar}}_{k,j}
\right)^2
\right|.
\end{equation}

\subsection{Total Normalized Loss}
The full loss is a sum of the spectral and local contributions. The implemented total loss is
\begin{equation}
\mathcal{L}
=
\frac{
\mathcal{L}_{\mathrm{deg}}
+
\mathcal{L}_{\mathrm{gap}}
+
\mathcal{L}_{Q}
+
\mathcal{L}_{\mathrm{MP}}
}{
\overline{|E|}+\eta
},
\end{equation}
where $\overline{|E|}$ is the mean absolute energy scale of the spectrum and $\eta$ is a small numerical offset. The spectral weights are already included in $\mathcal{L}_{\mathrm{deg}}$ and $\mathcal{L}_{\mathrm{gap}}$.

The normalization by $\overline{|E|}$ makes the spectral part of the loss less sensitive to the overall energy scale of a trial Hamiltonian. However, this normalization is not invariant under a uniform energy shift $H\mapsto H+\alpha\mathbb{1}$, since such a shift changes the absolute energies but leaves gaps and eigenvectors unchanged. In the parametrization used here there is no independent identity-shift parameter, so the optimizer does not explore such shifts directly.

A further limitation is that the same denominator also rescales the charge and Majorana-polarization penalties. These terms are dimensionless eigenvector diagnostics rather than energy penalties, so dividing them by an energy scale can make their contribution smaller when the spectrum becomes larger, even if the local Majorana quality does not improve. This possible bias is partly mitigated by fixing the hopping scale to $t=1$ and fixing $U$ during each optimization run, but the scalar loss should still be interpreted together with the separate spectral and local diagnostics.



\subsection{Search Strategy}
The resulting loss landscape is non-convex, especially once interactions are included. The numerical search therefore combines local gradient-based optimization with a basin-hopping procedure. The local optimizer is Adam, implemented with PyTorch automatic differentiation \cite{paszke2017automatic,kingma2014adam}.

At each Adam step, the reduced vector $\theta$ is converted into physical parameters, the Hamiltonian is constructed and diagonalized, and PyTorch automatic differentiation computes the gradient of the loss with respect to $\theta$.

The Adam update can be written as
\begin{equation}
\theta_k
=
\theta_{k-1}
-
\eta_{\mathrm{Adam}}
\frac{\hat{m}_k}{\sqrt{\hat{v}_k}+\varepsilon_{\mathrm{Adam}}},
\end{equation}
where $\eta_{\mathrm{Adam}}$ is the learning rate, and $\hat{m}_k$ and $\hat{v}_k$ are the bias-corrected moving averages of the gradient and squared gradient, respectively \cite{kingma2014adam}. In the implementation, Adam is called with learning rate $\eta_{\mathrm{Adam}}=0.05$, while the remaining Adam parameters are left at the PyTorch defaults, $\beta_1=0.9$, $\beta_2=0.999$, and $\varepsilon_{\mathrm{Adam}}=10^{-8}$. The parameters $\beta_1$ and $\beta_2$ set the moving averages that enter $\hat{m}_k$ and $\hat{v}_k$.

The numerical implementation of the Adam and basin-hopping search starts from a current parameter vector $\theta_{\mathrm{current}}$. A basin-hop is generated by drawing a trial vector
\begin{equation}
\theta_{\mathrm{trial}}
=
\theta_{\mathrm{current}}+0.25\,\xi,
\end{equation}
where the components of $\xi$ are normally distributed with zero mean and unit variance. Around this trial vector, we perform $6$ restarted Adam runs, each with an additional random perturbation,
\begin{equation}
\theta_{\mathrm{trial}}+\delta^{(r)},
\qquad
\delta_i^{(r)}\sim \mathrm{Unif}(-0.8,0.8).
\end{equation}
The scale $0.8$ was kept fixed as a practical perturbation scale, so that each basin-hop proposal is tested by several local Adam searches rather than a single trajectory. Each restarted run uses $300$ Adam iterations. The lowest-loss result from these restarts is used as the candidate for that basin-hopping step.

If the candidate has lower loss than the current vector, it is accepted. A candidate with higher loss is accepted with probability
\begin{equation}
p
=
\exp\left[
-\left(
\mathcal{L}_{\mathrm{cand}}-\mathcal{L}_{\mathrm{current}}
\right)
\right].
\end{equation}
This allows the search to occasionally leave a local basin instead of always moving downhill. The full basin-hopping search uses $30$ such steps for each pair $(n,U)$.

These search hyperparameters were not optimized systematically. They were kept fixed as practical settings chosen to balance exploration of the loss landscape against computational cost. For each pair $(n,U)$, the best configuration found by this procedure is stored together with its loss, parameter configuration, and reconstructed physical parameters.


\section{Post-Optimization Diagnostics}
The optimization procedure returns candidate parameter sets ranked by a scalar loss. Since this loss combines several criteria, a low value does not by itself show which Majorana signatures are actually present. For each pair $(n,U)$, the lowest-loss parameter set is therefore used to rebuild and diagonalize $H_{\mathrm{PMM}}(\theta_{\mathrm{opt}})$, after which the parity, charge, and Majorana-polarization diagnostics introduced in the theory chapter and used in the loss function above are evaluated separately.

For plotting and comparison, the parity-resolved spectrum is shifted so that the lowest state in either parity sector has energy zero. The reported spectral quantities are the splitting of the lowest even--odd pair,
\begin{equation}
\delta_0 =
\left|
E_0^{(e)}-E_0^{(o)}
\right|.
\label{eq:method_ground_state_splitting}
\end{equation}
The corresponding gap from this lowest pair to the nearest excited state is
\begin{equation}
\Delta_{\mathrm{gap}}
=
\min\left(E_1^{(e)},E_1^{(o)}\right)
-
\max\left(E_0^{(e)},E_0^{(o)}\right),
\label{eq:method_low_energy_gap}
\end{equation}
using the shifted energies. The local charge diagnostic is evaluated for the same lowest parity pair. The site-resolved charge difference is
\begin{equation}
q_i =
\left|
\bra{e_0}n_i\ket{e_0}
-
\bra{o_0}n_i\ket{o_0}
\right|.
\end{equation}
The summed charge difference reported in the results is
\begin{equation}
Q_0=\sum_{i=1}^{n} q_i .
\label{eq:method_charge_difference}
\end{equation}
In the diagnostic figures, the individual $q_i$ values are plotted together with the Majorana-polarization profile so that charge visibility and spatial localization can be compared site by site. From the lowest-pair Majorana-polarization profile $M_i$, the reported edge and bulk summaries are
\begin{equation}
\overline{|M_{\mathrm{edge}}|}
=
\frac{1}{2}
\left(
|M_1|+|M_n|
\right),
\label{eq:method_edge_polarization}
\end{equation}
and
\begin{equation}
\max |M_{\mathrm{bulk}}|
=
\max_{2\leq i\leq n-1}|M_i|.
\label{eq:method_bulk_polarization}
\end{equation}
For two-dot systems there are no interior sites, so the bulk contribution is set to zero.

The post-optimization diagnostics determine which optimized chains are used in the braiding calculations. Since the later sector-resolved braiding analysis tests projected manifolds beyond the lowest parity pair, the selected configurations are not judged only by the ground-state splitting, but also by the excitation gap, charge difference, and Majorana-polarization profile across the relevant low-energy structure.

\section{Effective Majorana Operator Construction}
\label{sec:method-effective-majorana-construction}
% 3.4 Effective Majorana Operator Construction
%   - energy-level Majoranas from even/odd pairs
%   - number of parity pairs included
%   - projected plus/minus components
%   - embedding into three-subsystem Hilbert space with JW strings
%   - normalization and algebra checks
The braiding calculations require explicit matrix representations of the Majorana operators associated with the optimized PMM chains. The theory chapter described why these operators should connect even and odd parity states. Here the focus is narrower: how the operators are constructed from the optimized eigenstates, embedded into the three-subsystem Hilbert space, projected to a chosen low-energy space, and checked before they are used in the braid.

\subsection{Energy-Level Majoranas from Parity Pairs}
For each optimized subsystem, the Hamiltonian is diagonalized and the eigenstates are separated into even and odd parity sectors. The states are then paired by their order in energy and inserted into the energy-level construction of Eq.~\ref{eq:majorana_construction}. In the implementation, the selected number of parity pairs is denoted by $m$, so the effective Majorana operators are built as
\begin{equation}
\gamma_{+}^{(m)}
=
\sum_{j=0}^{m-1}\gamma_{+,j},
\qquad
\gamma_{-}^{(m)}
=
\sum_{j=0}^{m-1}\gamma_{-,j}.
\label{eq:method_energy_level_majoranas}
\end{equation}
Here $\gamma_{\pm,j}$ denotes the contribution from the $j$th even--odd pair. In the code, $m$ is controlled separately from the total projection dimension: the projection dimension determines how much of the full three-subsystem Hilbert space is retained, while $m$ determines how many subsystem parity pairs are used to build the Majorana operator itself.

\subsection{Embedding Into the Three-Subsystem Hilbert Space}
Using the Jordan--Wigner ordering specified in Sec.~\ref{sec:method-jordan-wigner-construction}, subsystem operators are embedded into the full three-subsystem Hilbert space as
\begin{equation}
\gamma_A \mapsto \gamma_A\otimes I_B\otimes I_C,
\end{equation}
\begin{equation}
\gamma_B \mapsto JW_A\otimes \gamma_B\otimes I_C,
\end{equation}
and
\begin{equation}
\gamma_C \mapsto JW_{AB}\otimes \gamma_C,
\end{equation}
where $JW_A$ is the Jordan-Wigner string over subsystem $A$, and $JW_{AB}$ is the corresponding string over subsystems $A$ and $B$. The strings preserve the fermionic anticommutation relations between subsystems.

\subsection{Projection to the Retained Space}
The full three-subsystem Hamiltonian is diagonalized, and a projection basis $V$ is formed from a selected set of eigenvectors. In the low-energy benchmark, $V$ contains the eigenvectors of the retained lowest-energy manifold. In the sector-resolved scans, a separate basis $V_i$ is instead formed from all eigenvectors belonging to energy sector $i$. Here, sector $i$ denotes the $i$-th degenerate energy manifold, meaning the set of eigenvectors with the same energy within numerical tolerance. Any full-space operator $O$ is then represented in the selected space by
\begin{equation}
O_P = V^\dagger O V .
\label{eq:method_operator_projection}
\end{equation}
This projection is applied to the embedded Majorana operators before they are used in the projected braid models. The same projection basis is also used for all other operators in a given calculation, so that the energy-level Majoranas, local operators, static Hamiltonian, and driven braid terms act in the same retained space.

\subsection{Projected Local Components}
In addition to the energy-level construction above, the code constructs projected local Majorana components from the microscopic endpoint operators introduced in Eq.~\ref{eq:local_gamma_ops}. For a chosen projection basis $V$, the generic projected components are
\begin{equation}
\Gamma_{+,P}=V^\dagger(c^\dagger+c)V,
\qquad
\Gamma_{-,P}=V^\dagger i(c^\dagger-c)V .
\end{equation}
These projected local components provide a more microscopic reference for the Majorana operator at a chosen end of a subsystem. The direct projected local-operator model uses the normalized minus component $\Gamma_{-,P}$. When this operator is compared with the corresponding energy-level Majorana, only an overall sign is allowed,
\begin{equation}
s_{\mathrm{opt}}
=
\underset{s\in\{-1,+1\}}{\operatorname{argmin}}
\left\|
\Gamma_{-,P}-s\gamma_{-}^{P}
\right\|_F .
\end{equation}
This comparison fixes the arbitrary overall sign of the Majorana operator. It does not fit a general linear combination of \(\gamma_{+,P}\) and \(\gamma_{-,P}\).

For higher-dimensional energy sectors, a single overall sign became insufficient. Degenerate sectors contain several independent blocks associated with different combinations of subsystem energy-level labels. A label operator is therefore constructed from the paired eigenstates of each subsystem and projected into the selected sector. Diagonalizing this projected label operator decomposes the sector into labelled blocks. Within each block $r$, the projected local operator is fitted independently to the two available energy-level Majorana components,
\begin{equation}
\Gamma_{-,P}^{(r)}
\approx
a_r\gamma_{+,P}^{(r)}
+
b_r\gamma_{-,P}^{(r)} .
\label{eq:method_block_matched_majorana}
\end{equation}
The coefficients $a_r$ and $b_r$ are continuous least-squares coefficients and are allowed to differ between blocks. The fitted blocks are then reassembled into a block-matched local operator on the full selected sector. This construction tests whether the local operator can be represented by the energy-level Majorana pair after the independent degenerate blocks have been identified. The detailed construction of the label operator and the resulting block structure is presented and explained in the appendix \ref{app:block_label_method}.

\subsection{Normalization and Algebra Checks}
After projection into the retained subspace, the Majorana operators used in the braid are checked numerically. This includes the projected energy-level Majoranas and, when used, the projected local or block-matched operators. The diagnostics test Hermiticity, normalization,
\begin{equation}
\gamma_i^2 \approx I,
\end{equation}
and mutual anticommutation,
\begin{equation}
\{\gamma_i,\gamma_j\}\approx 0
\qquad (i\neq j).
\end{equation}
These algebra checks quantify how well the projected operators retain the Majorana structure in the chosen subspace.


\section{Braiding Models and Pulse Protocol}

The braiding calculations use the optimized PMM chains as building blocks for a three-subsystem junction. The purpose of this section is to define the Hamiltonians that are driven during the braid and the common pulse sequence used for all model variants. The following section then describes how the resulting time-dependent evolution is calculated and compared with the target exchange gate.

\subsection{Three-Subsystem Geometry}
The junction consists of three optimized PMM subsystems, labelled $A$, $B$, and $C$. For each subsystem, energy-level Majorana operators are constructed from the even--odd parity pairs as described in Sec.~\ref{sec:method-effective-majorana-construction}. In the braid notation used below, the two energy-level Majoranas from subsystem $A$ are denoted by $\gamma_0$ and $\gamma_1$. The Majoranas associated with subsystems $B$ and $C$ are denoted by $\gamma_2$ and $\gamma_3$ when they are used as the exchanged pair.


In this notation, the intended exchange is generated by moving the coupling of $\gamma_0$ from $\gamma_1$ to $\gamma_2$, then to $\gamma_3$, and finally back to $\gamma_1$.

All braiding models are built in the projected Hilbert space introduced in Sec.~\ref{sec:method-effective-majorana-construction}.

The driven braid terms are supplemented by the projected static Hamiltonian of the uncoupled subsystem background,
\begin{equation}
H_{\mathrm{static}}
=
V^\dagger\left(H_B+H_C\right)V.
\end{equation}
This term keeps the internal optimized-chain structure of subsystems $B$ and $C$ while the time-dependent junction couplings are varied. Since it is included in the same way for the model comparisons, the differences between the braid results come from the choice of driven coupling operators rather than from a different static background.

\subsection{Energy-Level Reference Model}
The reference braid model uses four energy-level Majorana operators and hand-chosen bilinear couplings. In the projected space, the time-dependent Hamiltonian is
\begin{equation}
H_{\mathrm{ref}}(t)
=
\lambda_1(t)\, i\gamma_0\gamma_1
+
\lambda_2(t)\, i\gamma_0\gamma_2
+
\lambda_3(t)\, i\gamma_0\gamma_3 .
\label{eq:method_ideal_braid_hamiltonian}
\end{equation}
This model is used as a reference calculation. Since the couplings are written directly as Majorana bilinears, it isolates the braid protocol itself from the additional complications of microscopic junction operators and imperfect projected Majoranas.

\subsection{Projected Microscopic Junction Model}
The projected microscopic junction model starts from the physical endpoint
couplings between the optimized PMM subsystems. We use the junction operator
introduced in Eq.~\ref{eq:junctionHamiltonian}, with an analogous operator for
the $A$--$C$ junction, and project these full-space operators into the retained
low-energy space,
\begin{equation}
T_{AB}=V^\dagger H^{\mathrm{junc}}_{AB}V,
\qquad
T_{AC}=V^\dagger H^{\mathrm{junc}}_{AC}V .
\label{eq:method_projected_junctions}
\end{equation}

These two terms can help identify which low-energy Majorana channel is selected by the physical endpoint coupling. We can decompose the projected junction operators $T_{AB}$ and $T_{AC}$, into the basis of Majorana bilinears,
\begin{equation}
  c_{\alpha}^{AB} = \frac{1}{d} \Tr\left[(i\gamma_{A1}\gamma_{B\alpha})^{\dagger} T_{AB}\right],
    \qquad
    c_{\beta}^{AC} = \frac{1}{d} \Tr\left[(i\gamma_{A1}\gamma_{C\beta})^{\dagger} T_{AC}\right],
\end{equation}
testing for both $\alpha,\beta \in \{1,2\}$, where $d$ is the dimension of the projected space. The normalized trace is the Hilbert--Schmidt overlap between the projected junction operator and a Majorana bilinear\cite{hazewinkelEncyclopaediaMathematics2002}. The $\gamma_{S,i}$ operators are the projected energy-level Majoranas from the subsystems $S=A,B,C$, constructed from the parity-pair operators of Eq.~\ref{eq:majorana_construction} as summarized in Eq.~\ref{eq:method_energy_level_majoranas}. The maximum absolute value of the coefficients $c_{\alpha}^{AB}$ and $c_{\beta}^{AC}$ identifies which low-energy Majorana channel is selected by the physical endpoint coupling.


The projected microscopic braid Hamiltonian is then
\begin{equation}
H_{\mathrm{junc}}(t)
=
\lambda_1(t)\,T_A
+
\lambda_2(t)\,T_{AB}
+
\lambda_3(t)\,T_{AC},
\label{eq:brajunc}
\end{equation}
where $T_A=i\gamma_0\gamma_1$ is the initial coupling inside subsystem $A$.
This model keeps the full projected hopping-and-pairing structure of the
microscopic junctions.

\subsection{Projected Local-Operator Model}
The projected local-operator model uses the channel identified by the projected
microscopic junction decomposition, but constructs the exchanged operators from
local endpoint Majorana components instead of from the full junction
Hamiltonian. The endpoint components are those defined in
Eq.~\ref{eq:local_gamma_ops}. After projection and normalization, the component
selected by the junction analysis defines the local Majorana operator
$\Gamma_S^P$ used in the braid.

The corresponding local-operator braid uses the same bilinear form as
Eq.~\ref{eq:local_braid_ham}. In the method notation, this corresponds to the
identifications $\gamma_{A1}\rightarrow\gamma_0$,
$\gamma_{A2}\rightarrow\gamma_1$, $\lambda_A\rightarrow\lambda_1$,
$\lambda_{AB}\rightarrow\lambda_2$, and
$\lambda_{AC}\rightarrow\lambda_3$.
This construction tests whether the same exchange channel selected by the full
projected junction is also obtained from projected local endpoint Majorana
operators.


\subsection{Block-Matched Local-Operator Model}
The sector-resolved scans also use a block-matched local-operator model. It has the same bilinear form as Eq.~\ref{eq:local_braid_ham}, but replaces the direct projected local operators by the block-matched operators defined in Eq.~\ref{eq:method_block_matched_majorana}. The labelled blocks are fitted independently and then reassembled before the braid Hamiltonian is constructed. This model tests whether resolving the internal degeneracy structure of a sector improves the braid generated from the local channel.

\subsection{Pulse Sequence}
The same pulse sequence is used for each braid model. To avoid confusing the braid amplitudes with the superconducting pairing parameters $\Delta_i$, the time-dependent braid couplings are denoted by $\lambda_i(t)$. The coupling envelope is a smooth pulse constructed from two sigmoid functions,
\begin{equation}
\lambda(t;T_p)
=
\lambda_{\min}
+
(\lambda_{\max}-\lambda_{\min})
\frac{1}{1+\exp[-sP_+(t;T_p)]}
\frac{1}{1+\exp[sP_-(t;T_p)]},
\label{eq:method_lambda_pulse}
\end{equation}
where
\begin{equation}
P_+(t;T_p) = t-T_p+w/2,
\qquad
P_-(t;T_p) = t-T_p-w/2.
\end{equation}
Here $T_p$ is the pulse center, $w$ is the pulse width, and $s$ controls the steepness of the rise and fall. The three braid couplings are then chosen as
\begin{equation}
\lambda_1(t)
=
\lambda(t;0)+\lambda(t;T)-\lambda_{\min},
\qquad
\lambda_2(t)
=
\lambda(t;T/3),
\qquad
\lambda_3(t)
=
\lambda(t;2T/3).
\label{eq:method_braid_pulses}
\end{equation}
Thus the $A$-arm coupling is active at the beginning and end of the cycle, while the $A$--$B$ and $A$--$C$ couplings are activated in sequence during the middle of the protocol. In the numerical calculations used here, the total braid time is set to $T=1$, with $\lambda_{\max}=1$, $\lambda_{\min}=0$, pulse width $w=T/3$, and steepness $s=20/w$.

At each time step, the Hamiltonian for the chosen model is assembled from the same three time-dependent coefficients. The instantaneous spectrum is monitored along the path, and the isolated low-energy band is then transported using the Kato procedure described in the next section.

\section{Kato Transport and Braid Validation}

After a braid Hamiltonian has been defined, the remaining task is to compute the unitary generated by the adiabatic path and to check whether it implements the desired Majorana exchange. The calculation is performed in the projected Hilbert space, but the transported subspace can be smaller than the full projection space. This section describes the transported bands used in the calculations, how the Kato evolution is computed, and which diagnostics are used to compare the resulting unitary with the target braid.

\subsection{Transported Bands}
In the low-energy benchmark, the transported subspace is the isolated four-dimensional lowest-energy manifold. Its gap to the remaining projected states is monitored throughout the braid path.

In the sector-resolved scans, each energy sector is projected and tested independently. Activating the braid Hamiltonian splits the initially degenerate sector into a lower and an upper band of equal dimension. The Kato calculation transports the lower band, so that for a sector projection of dimension $d_{P_i}$ the transported dimension is fixed to
\begin{equation}
d_{T_i}=\frac{d_{P_i}}{2}.
\label{eq:method_sector_transport_dimension}
\end{equation}
The spectral gap between these two bands is monitored along the path. The existence of this isolated lower band makes the transport well defined, while the final target comparison determines whether the transported unitary actually performs the intended exchange.

\subsection{Numerical Kato Transport}
The braid path is discretized into time points $t_0,\ldots,t_N$. At each time step, the Hamiltonian is diagonalized and the projector onto the transported band is constructed,
\begin{equation}
P(t_\ell)
=
\sum_{a=1}^{d_T}
\ket{\psi_a(t_\ell)}\bra{\psi_a(t_\ell)} .
\end{equation}
The Kato generator is approximated from consecutive projectors,
\begin{equation}
K(t_\ell)
=
P(t_\ell)\dot{P}(t_\ell)
-
\dot{P}(t_\ell)P(t_\ell),
\qquad
\dot{P}(t_\ell)
\approx
\frac{P(t_{\ell+1})-P(t_\ell)}{\Delta t}.
\end{equation}
The transport unitary is updated as
\begin{equation}
U_K(t_{\ell+1})
=
\exp[-\Delta t\,K(t_\ell)]U_K(t_\ell),
\qquad
U_K(t_0)=I .
\label{eq:method_kato_update}
\end{equation}
This gives the geometric transport of the selected subspace along the braid path. The instantaneous energies and pulse amplitudes are stored at the same time, so that the quality of the transported band can be checked together with the final braid unitary.

\subsection{Target Exchange Gate}
The target unitary for exchanging $\gamma_2$ and $\gamma_3$ is taken to be
\begin{equation}
U_{\mathrm{target}}
=
\exp\left(
-\frac{\pi}{4}\gamma_2\gamma_3
\right).
\label{eq:method_target_exchange_gate}
\end{equation}
This convention corresponds to the exchange action
\begin{equation}
\gamma_2\mapsto -\gamma_3,
\qquad
\gamma_3\mapsto \gamma_2,
\end{equation}
up to the overall sign convention chosen for the Majorana operators. This is the Heisenberg action $U_K^\dagger\gamma_i U_K$ used in the numerical checks, which is the inverse conjugation convention to the one written in Sec.~\ref{sec:theory-majorana-exchange-gates}. Since a global phase of the transported unitary is physically irrelevant, comparisons with the target gate are made after aligning the overall phase. In practice, the reported gate error is a Frobenius-norm distance between the transported unitary and the phase-aligned target unitary, restricted to the initial transported basis.

For the sector-resolved scans, the agreement with the target is also summarized by the normalized overlap
\begin{equation}
\mathcal{O}
=
\frac{1}{d_T}
\left|
\operatorname{Tr}\left(U_{\mathrm{target}}^{\dagger}U_K\right)
\right|,
\qquad
D_{\mathcal{O}}=|1-\mathcal{O}|,
\label{eq:method_overlap_deviation}
\end{equation}
where \(d_T\) is the dimension of the transported subspace. A perfect match up to a global phase gives \(\mathcal{O}=1\), so the plotted quantity \(D_{\mathcal{O}}\) vanishes for a perfect target match and increases as the transported unitary deviates from the target.

\subsection{Operator-Action Checks}
The first validation step checks the induced action on the Majorana operators directly. For each model, the transported unitary is applied as
\begin{equation}
\gamma_i'
=
U_K^\dagger \gamma_i U_K.
\end{equation}
The resulting operators are compared with the expected single-exchange map,
\begin{equation}
\gamma_2' \approx -\gamma_3,
\qquad
\gamma_3' \approx \gamma_2,
\qquad
\gamma_0' \approx \gamma_0,
\qquad
\gamma_1' \approx \gamma_1.
\end{equation}
The mismatch is measured as the normalized Frobenius norm of each operator-map residual. For the exchanged Majoranas,
\begin{equation}
\epsilon_{\gamma_2}
=
\frac{
\left\|U_K^\dagger \gamma_2 U_K+\gamma_3\right\|_F
}{\sqrt{d}},
\qquad
\epsilon_{\gamma_3}
=
\frac{
\left\|U_K^\dagger \gamma_3 U_K-\gamma_2\right\|_F
}{\sqrt{d}} .
\end{equation}
The spectator Majoranas are checked in the same way,
\begin{equation}
\epsilon_{\gamma_0}
=
\frac{
\left\|U_K^\dagger \gamma_0 U_K-\gamma_0\right\|_F
}{\sqrt{d}},
\qquad
\epsilon_{\gamma_1}
=
\frac{
\left\|U_K^\dagger \gamma_1 U_K-\gamma_1\right\|_F
}{\sqrt{d}} .
\end{equation}
Here $d$ is the matrix dimension in the comparison space.
These quantities report the braid error directly in the Majorana-operator language used to define the target exchange.
In the projected calculations, the same operator-action check is evaluated with the target Majoranas appropriate to the construction being tested. For the energy-level construction, $\gamma_2$ and $\gamma_3$ are the Majorana operators constructed from even--odd parity pairs, as described in Sec.~\ref{sec:method-effective-majorana-construction}. For the projected-local and block-matched constructions, the corresponding target operators are instead built from the projected local endpoint Majorana components.

The mismatch between the energy-level target and the projected-local target is also recorded. A small mismatch indicates that the two constructions define the same effective exchange channel in the retained subspace, while a larger mismatch means that the target definition itself contributes to the reported braid error.

\subsection{Parity-Resolved Gate Checks}
The transported unitary is also analyzed in a parity basis. The total parity operator is projected into the initial transported space and diagonalized, separating the unitary into odd- and even-parity blocks. The off-block leakage is defined from the norm of the odd-to-even and even-to-odd blocks,
\begin{equation}
\mathcal{L}_{\mathrm{parity}}
=
\|U_{oe}\|
+
\|U_{eo}\|.
\end{equation}
A small value indicates that the braid approximately preserves fermion parity inside the transported manifold. This is required by fermion-parity conservation and therefore serves as a numerical check.

After the parity blocks have been separated, each block is compared with the corresponding block of the target exchange gate in Eq.~\ref{eq:method_target_exchange_gate}. This produces separate odd- and even-sector target-gate errors. Reporting the parity blocks separately is useful because the same Majorana exchange can appear with different matrix representations in the two parity sectors, while still describing the same physical braid.
